\chapter*{Introduction}
\addcontentsline{toc}{chapter}{Introduction}
The theoretical framework of Supersymmetry is based on the existence of a 
new kind of symmetry in which to each bosonic particle
corresponds a new fermion and to each fermionic particle corresponds a new boson:
these are called superpartner particles and their spin differs from their
"regular" partner's by a half integer. More precisely, Supersymmetry was
introduced in the context of Quantum Field Theories as an internal symmetry
of a Lagrangian between fermionic and bosonic degrees of freedom, as a way to
unify interactions (bosons) with matter (fermions).\\
So far in nature none of these superpartners have been observed and so
it is conjectured that Supersymmetry must be spontaneously broken and
the scale of this breaking is at least of the order of $100$ GeV.

From this context, researchers started to test supersymmetric theories in
simpler theoretical frameworks, like that of non-relativistic Quantum
Mechanics, finding insight on solvable potentials.

In the first chapter of this thesis the formalism and algebraic structure of
Supersymmetric Quantum Mechanics is introduced and the properties of
supersymmetric partner Hamiltonians are stated and proven; some examples
are then provided and explicitely explored. The main references for this
chapter are \cite{cooper_main} and \cite{khare_short} for the definitions and
properties of Supersymmetry and \cite{bastianelli} for the Supersymmetric
Algebra.

In the second chapter of this thesis the concept of shape invariant partner
potentials is introduced with their properties: this definition is then used
to construct a generalization of the ladder operators for the Quantum Harmonic
Oscillator to a greater set of integrable 1-dimensional potentials; finally, some
examples of potentials with complex analytical solutions are solved with the
properties of shape invariance. The main reference for this chapter is again
\cite{cooper_main}. 

In the third and final chapter of this thesis a variation
of the WKB quantization condition on bound states of supersymmetric 
partner potentials is derived for both broken and unbroken supersymmetry; it is
then shown that this method yields the exact energy eigenvalues of the spectrum
of shape invariant potentials. The references for this chapter are \cite{zwiebach}
for the general theory on the WKB method and \cite{cooper_main}, \cite{dutt_swkb}
for the supersymmetric approach.
